% Auto-generated by thesis/make_prompts.py --- do not edit by hand.
% Source: analyse_ofsted_reports.py, analyse_behaviour_policies_v3.py

Both scorers share a prompt architecture. A system instruction establishes
the task, the scoring procedure and the output contract; the substantive
content is delivered as a single JSON payload carrying the school
identifiers, the extracted document text, the full rubric, and a
specification of the required output fields. The model must return
schema-valid JSON --- through the provider's structured-output facility
where available, and through a forced tool call otherwise --- so that no
free-text parsing is involved. Sampling temperature is zero. Each request is
keyed by a SHA-256 hash of the model name, prompt version, schema version and
payload, so re-running a scorer over unchanged documents returns cached
scores rather than silently re-scoring them.

Both rubrics require the model to state, for every score, why the score is
not one point higher and not one point lower. Both apply the same borderline
rule: where evidence is genuinely ambiguous, the lower score is taken. The
asymmetric exception in both is that a score below the midpoint requires
positive evidence of a deficiency --- silence on behaviour, or on
relationships, is scored as unremarkable rather than as poor. This matters
for interpreting the warmth null in \cref{sec:p1_text_valid}: the rubrics are
constructed so that a document which simply does not discuss warmth cannot
produce a low warmth score, and the observed compression at the midpoint is
therefore the designed response to silence, not a scoring failure.

\subsection*{Ofsted reports}

Prompt version \texttt{ofsted-strictness-warmth-teaching-v3}. System instruction:

\begin{quote}\small
You are scoring Ofsted report evidence for research. Use only the supplied passages and rubric. Apply the hard rules exactly. You must call the record\_scores tool with your scores --- do not reply in plain text.
\end{quote}

The payload supplies the URN and school name, inspection date and report
type, Ofsted's own behaviour and attitudes judgement (as context only, and
explicitly not as a cap on the score), the opening summary, the
improvement-areas section, and the substantive inspection narrative with
boilerplate, primary-phase and sixth-form passages removed. For an
all-through school a further note instructs the model to score only the
secondary provision described in the remaining text. The required output is a
score, a confidence level and a reason for each of strictness, warmth and
teaching, plus a manual-review flag.

\paragraph{Strictness scale anchors}
\begin{description}
\item[Score 1] Rules are unclear, inconsistently applied, or disruption is a persistent and documented problem.
\item[Score 2] Standards exist but enforcement is patchy --- some year groups, subjects, or staff are exceptions; behaviour is an identified area for improvement. Note: Ofsted inspections take place when schools are typically presenting their best behaviour --- pupils and staff are aware they are being observed. Direct observations by inspectors of low-level disruption, inattention, or inconsistent enforcement during the visit therefore carry more evidential weight than they might appear to, and should weigh toward S=2 rather than S=3 even if described mildly.
\item[Score 3] Clear expectations are in place and generally maintained; the evidence mentions behaviour or conduct but without singling it out as a strength. Adequate, unremarkable enforcement.
\item[Score 4] Multiple evidence passages consistently indicate that conduct standards are above average --- the evidence as a whole suggests behaviour is a strength, with adults clearly driving high expectations. Does not require verbatim ``strength'' language but the pattern across passages must clearly point above adequate. If behaviour outside the classroom --- corridors, social times, around the building --- is explicitly flagged as an improvement area, this prevents a 4; above-average conduct must be consistent across the whole school environment, not just in lessons.
\item[Score 5] The evidence emphatically and consistently highlights conduct as exceptional across most or all passages --- the school's strictness is unmistakably a defining feature of its character.
\end{description}

\paragraph{Warmth scale anchors}
\begin{description}
\item[Score 1] Staff-pupil relationships are poor, transactional, or pupils feel unknown or uncared for by staff.
\item[Score 2] The report contains positive evidence of relational inconsistency: some pupils feel unknown to, disconnected from, or let down by staff. Even mild or qualified negative relational language --- pupils feeling their concerns are not always heard, that bullying is not resolved, or that staff are not accessible --- is sufficient for W=2 and is not offset by generic positive language about safety or pastoral provision.
\item[Score 3] The report describes a safe and generally positive environment. Positive language is present but it characterises the school's overall character or provision rather than the quality of actual staff-pupil interactions. W=3 language: the school described as ``calm'', ``caring'', ``warm'', ``welcoming'', or ``inclusive'' used as a general school-level label; pupils feeling safe, happy, settled, or supported stated as an outcome without describing how staff create that; generic statements that staff are supportive, approachable, or accessible (availability, not relational quality); PSHE, pastoral teams, heads of year, or trusted-adult arrangements mentioned or described positively; enrichment activities or school events. None of these describe relational quality --- they describe provision, character labels, or bare outcomes. Note: any credible negative relational signal --- pupils feeling bullying is not resolved, that they do not feel safe, or that staff are not available --- should pull the score toward W=2.
\item[Score 4] The evidence describes the quality of staff-pupil relationships in explicit terms that go beyond characterising the school's overall character or stating that systems exist. W=4 language describes HOW staff and pupils actually relate --- not just that a warm school exists or that pastoral provision is in place. Qualifying language includes: staff described as knowing pupils well; inspectors observing or describing warm, caring, or positive interactions between staff and pupils; strong, warm, or positive staff-pupil relationships described as a characteristic of how adults and pupils interact (not merely as a school-label adjective); pupils described as feeling genuinely valued, known, or cared for in ways that attribute this to staff behaviour or the relational quality rather than to provision existing; staff described as caring about pupils as individuals; pupils having positive or trusting relationships with staff. DOES NOT QUALIFY for W=4 (stays at W=3): ``caring school'', ``warm community'', ``inclusive ethos'', ``staff are approachable / accessible'', ``pupils can speak to a trusted adult'', ``pastoral team is effective'', ``bullying is dealt with swiftly'', ``PSHE is taught well'', ``there is a positive school community''. These are provision-existence or school-character statements, not descriptions of relational quality. No significant relational negatives are present. Where both positive and negative relational evidence appear, the score is W=3.
\item[Score 5] Relational warmth is a prominent and explicitly named feature of the school's character, evidenced across multiple passages. To reach W=5, warmth or the quality of staff-pupil relationships must be described as a standout, defining, or exceptional feature of the school --- not merely present or adequate. Markers: inspectors use language such as ``exceptionally caring'', ``deeply supportive'', ``truly know their pupils'', ``outstanding relationships'', or equivalent; multiple passages independently converge on the same positive relational theme; or warmth is explicitly named as one of the school's defining characteristics. A single warm passage alongside generic positives does not reach W=5. The absence of negatives alone is never sufficient --- positive evidence of distinctive relational quality must be present.
\end{description}

\paragraph{Scoring rules: strictness and warmth}
\begin{enumerate}
\item Use only the supplied evidence.
\item Do not infer beyond the supplied evidence.
\item The Ofsted Behaviour and attitudes judgement is provided as context only; do not treat it as a hard cap on your score.
\item BORDERLINE RULE: When evidence is genuinely ambiguous or limited, always assign the lower score. For scores that sit on the boundary between two adjacent values --- where reasonable arguments exist for both --- always choose the lower value. A higher score is warranted only when the evidence clearly and unambiguously supports it. This is a deliberate design choice for consistency. IMPORTANT EXCEPTION FOR STRICTNESS: scores of 1 or 2 each require positive evidence that behavioural problems exist (documented disruption, patchy enforcement, behaviour identified as an improvement area). If behaviour is simply not discussed in the report --- because the overall picture is positive and inspectors focused elsewhere --- do not score below 3. Absence of behaviour-focused writing is not evidence of poor behaviour; default to 3 (adequate, unremarkable) in that case. IMPORTANT EXCEPTION FOR WARMTH: scores of 1 or 2 each require positive evidence of poor or absent staff-pupil relationships --- explicit descriptions of pupils feeling uncared for, unknown to staff, or unsafe due to relational failures, or of staff-pupil relationships being poor or transactional. If warmth is simply not discussed in the report --- because inspectors focused on curriculum, outcomes, or other aspects --- do not score below 3. Absence of warmth-focused writing is not evidence of poor warmth; default to 3 (safe and generally positive, no direct relational description) in that case.
\item For each final score, explain briefly why the score is not one point higher and not one point lower.
\item PUPIL VIEWS ON RULES ARE WARMTH EVIDENCE ONLY: How pupils perceive, accept, or resent the school's rules and strictness must feed exclusively into the warmth score and must have no effect on the strictness score in any direction --- it must not raise it, lower it, or prevent it from reaching any threshold, including 5. Strictness is scored solely on what adults do: the standards they set, how consistently they enforce them, and the resulting behavioural outcomes. Whether pupils welcome or resent those standards is a relational question, not a strictness question. If you find yourself citing pupil views about rules as a reason to withhold a strictness score of 5, you are violating this rule --- remove that reasoning and re-evaluate based only on what the adults do.
\item WARMTH EXCLUSION (NARROW): SEND provision, knowledge of individual educational needs, mental health promotion programmes, and specialist provision for vulnerable groups are NOT evidence of warmth. General pastoral systems (e.g. trusted adults, pastoral teams described as accessible to all pupils), PSHE/PSHCE programmes, and whole-school care structures CAN be counted as warmth evidence --- but their existence qualifies only for W=3. Warmth is primarily evidenced by descriptions of the relational quality between staff and pupils: explicit kindness language, staff caring for pupils as individuals, pupils feeling genuinely valued and respected, gratitude practices, or direct statements that staff have pupils' best interests at heart.
\item W=3/W=4 THRESHOLD --- RELATIONAL QUALITY VS PROVISION OR CHARACTER LABEL: The most common scoring error is treating positive language about the school's general character or systems as evidence of relational warmth. The test is: does the evidence describe HOW staff and pupils actually relate --- or does it describe the school's general character, what provision exists, or aggregate outcomes without attributing them to relational quality? STAYS AT W=3: ``calm and caring school'' / ``warm and welcoming'' (school-level label); ``pupils feel safe and happy'' (bare outcome); ``staff are supportive / accessible / approachable'' (availability, not relational quality); ``bullying is dealt with effectively''; ``PSHE is in place''; ``pastoral team / heads of year are effective''; ``there is a positive school community''; ``pupils trust staff to resolve concerns''. QUALIFIES FOR W=4: ``staff know pupils well''; ``strong / warm / positive staff-pupil relationships'' (describing the relationship itself, not a school label); ``inspectors observed warm or caring interactions between staff and pupils''; ``pupils feel genuinely valued and known by staff''; ``staff care about pupils as individuals''; ``pupils have positive, trusting relationships with staff''; ``mutual respect between staff and pupils''. The test is whether the language characterises the relational quality itself --- not just the school's overall tone or what provision is available.
\end{enumerate}

\paragraph{Teaching scale anchors}
\begin{description}
\item[Score 1] Seriously weak teaching, directly harming learning. The report explicitly names teaching quality as a serious failure. Key markers: teacher subject knowledge is described as varying greatly or extremely weak in places; explanations and classroom delivery are described as failing to give pupils the knowledge they need (``work does not enable pupils to gain the knowledge needed for future success''); there is no effective checking of understanding. Outcome language attributes weak results directly to teaching failure (``pupils ill-prepared'', ``many do not receive an acceptable standard of education'', ``achievement is weak''). The picture across the report is one of teaching as a root cause of poor outcomes, not merely a symptom.
\item[Score 2] Inconsistent or developing teaching, named as requiring improvement. Teaching delivers the curriculum in some classrooms but not consistently. Key markers: 'some teachers deliver content clearly... this is not always the case'; checking for understanding is incomplete or ineffective (``do not always check before moving on'', ``feedback not as helpful as it could be'', ``gaps in pupils'' knowledge not addressed swiftly'); teacher adaptation is named as inconsistent (``some staff do not adjust their teaching in response to checks''). Teaching quality is named explicitly as an area requiring improvement. Outcome language describes pupils as ``not achieving as well as they should''. The report may acknowledge pockets of effective practice but the prevailing picture is inconsistency.
\item[Score 3] Adequate and broadly effective. DEFAULT score. Teaching is broadly adequate: the curriculum is delivered competently, teachers have secure subject knowledge, some checking for understanding takes place. Teaching quality is neither praised as a strength nor named as a weakness --- it is either unremarkable or discussed only in terms of curriculum design and sequencing rather than classroom delivery. Pupils make broadly expected progress. DEFAULT: If the report does not specifically discuss how teachers teach in the classroom --- because inspectors focused on curriculum intent, leadership, or behaviour --- assign score 3. Absence of teaching-focused writing is not evidence of poor teaching.
\item[Score 4] Strong classroom teaching, described positively across most of the school. The report describes teaching positively and with some specificity. Key markers: teachers have ``strong subject knowledge'' and ``know how best to share it''; checking for understanding is consistent (``teachers systematically check pupils'' understanding', ``misconceptions are quickly addressed''); teaching is effective for most pupils including those with SEND (``teachers meet needs securely in lessons''); explanations are described as clear and effective. The word ``typically'' or ``most teachers'' appears alongside positive teaching descriptions. Outcomes are good. The report may note isolated inconsistencies but the overall picture of classroom delivery is clearly positive.
\item[Score 5] Expert teaching, consistently described as a school strength. Teaching is described as expert or excellent and as a defining quality of the school. Key markers: ``expert teachers'', ``excellent knowledge''; language of precision and responsiveness (``explain complex ideas very well'', ``deliver subject content with precision'', ``nimble and responsive in adjusting teaching'', ``intervene quickly to address any misconceptions''); formative assessment is systematic and sharp (``teachers check understanding regularly'', ``assessment sharply identifies areas for improvement'', ``precisely tailored adaptations''); ``consistently'' or ``across all subjects'' accompanies praise for teaching. Outcome language explicitly and emphatically attributes strong results to teaching quality. The report leaves no doubt that teaching is an exceptional strength.
\end{description}

\paragraph{Scoring rules: teaching}
\begin{enumerate}
\item Use only the supplied evidence. Do not infer beyond what the report states.
\item DELIVERY vs DESIGN: Curriculum design language (intent, sequencing, what is taught) is background --- not itself evidence of teaching quality. Score only on language about HOW teachers teach in classrooms: subject knowledge, explanations, checking for understanding, adaptation to pupils, and response to gaps or misconceptions. Exception: ``curriculum delivery'' statements that describe what teachers actually do in practice (e.g. ``teachers deliver subject content with precision'') do count as delivery evidence.
\item CONSISTENCY LANGUAGE: Weigh consistency modifiers heavily. ``Consistently / routinely / across all subjects'' alongside positive language $\rightarrow$ push toward 5. ``Typically / most teachers'' with positive language $\rightarrow$ push toward 4. ``Some teachers / not always / inconsistent'' $\rightarrow$ push toward 2-3. ``Many do not / often ineffective'' $\rightarrow$ push toward 1-2.
\item CHECKING FOR UNDERSTANDING is the most reliably graded marker in Ofsted reports. If checking is described as sharp, precise, and responsive $\rightarrow$ supports 4-5. If described as present but incomplete or slow $\rightarrow$ supports 3-4. If described as absent, ineffective, or not acted on $\rightarrow$ supports 1-2.
\item BORDERLINE RULE: When delivery evidence is genuinely limited or absent, default to 3. When evidence points between two scores, prefer the lower unless outcome language explicitly and emphatically credits teaching quality.
\item For the final teaching score, explain briefly why it is not one point higher and not one point lower.
\end{enumerate}

\subsection*{Behaviour policy documents}

Prompt version \texttt{behaviour-policy-strictness-warmth-v26}. The system instruction
carries the scoring procedure itself, as an explicit decision tree rather than
as holistic guidance:

\begin{quote}\small
You are scoring school behaviour policy evidence for academic research. Use the anchor examples in the rubric as calibration references. Score strictness and warmth independently.

STRICTNESS DECISION PROCEDURE --- follow these steps in order:

Step 1: Does the policy describe a day-to-day detention or consequence system? (Not merely mentioning detentions exist; not merely listing what offences lead to suspension or permanent exclusion.) If NO --- score S=1 or S=2. S=1 = no consequence language at all (empty policy, rewards-only, purely aspirational with no sanctions content). S=2 = some consequence language is present but no coherent day-to-day system is described --- this includes policies with vague/proportionate sanctions language, restorative-primary approaches, or exclusions-only content.

Step 2 (only if Step 1 = YES): Is there a named tracking system (demerits, BfL points) with explicit numeric thresholds that automatically generate a named consequence? If NO --- score S=3. If YES --- go to Step 3.

Step 3 (only if Step 2 = YES): Does the demerit trigger list include ANY of: posture or seating position; eye contact/tracking/SLANTing; written work neatness or quality; corridor conduct (wrong side, running); minor uniform details (untucked shirt, incorrect shoe)? Note: qualifiers such as ``persistently'', ``repeatedly'', or ``sloppy'' specify the INFRACTION THRESHOLD, not whether staff choose to sanction --- they do NOT introduce discretion. If YES --- score S=5. If NO --- score S=4.

WARMTH DECISION PROCEDURE --- follow these steps in order:

Step 1: Does the policy contain positive culture content BEYOND aspirational sentences that describe the school's VALUES or EMPHASIS without naming any mechanism? The following do NOT count regardless of how elaborate they are: ``Emphasis is on recognising and celebrating effort and success so that all members feel valued and rewarded''; ``We develop all pupils in the habits of self-discipline and kindness''; ``students who behave well are likely to receive a reward''; any sentence describing what the school VALUES about recognition, success, or positive culture without naming specific rewards, pastoral structures, or staff practices. The following DO count: named rewards (house points, certificates, commendation events); named pastoral services (counselling, SEMH, safe spaces, pastoral programmes); explicit HOW-to staff guidance (operational, not aspirational). If NO qualifying content $\rightarrow$ W=1 [STOP --- do not proceed further]. If YES $\rightarrow$ Step 2.

Step 2a --- REWARDS CHECK: Does the policy describe a NAMED rewards mechanism with named types (house points, certificates, commendation events, reward tokens) AND described occasions? If YES $\rightarrow$ Step 2b. If NO $\rightarrow$ Step 2c.

Step 2b (rewards present) --- INDEPENDENT PASTORAL CHECK: Are there also INDEPENDENTLY ACCESSIBLE pastoral/wellbeing structures --- counselling service, SEMH lead, wellbeing programme, safe space available to ALL students independent of sanction status, dedicated pastoral programme? NON-QUALIFYING (consequence-integrated only): HoY/SLT reports for sanctioned students, homework clubs for non-compliant students, reintegration meetings after removal, restorative conversations within the consequence procedure, self-control mentoring triggered by poor behaviour, support strategies mentioned in exclusion procedure text. If YES independent pastoral $\rightarrow$ Step 3. If NO independent pastoral $\rightarrow$ W=3.

Step 2c (NO rewards) --- INDEPENDENT PASTORAL CHECK: Are there INDEPENDENTLY ACCESSIBLE pastoral/wellbeing structures (counselling, SEMH, wellbeing programme, non-consequence-triggered safe spaces --- as defined above)? If YES $\rightarrow$ W=3 [STOP --- do not proceed to Step 3]. If NO $\rightarrow$ W=2 [STOP --- do not proceed to Step 3]. (Consequence-integrated supports alone, however numerous, cannot exceed W=2.)

Step 3: Is there explicit operational guidance on HOW staff should relate to individual students --- instructional (``learn their names'', ``use humour --- it builds bridges'', ``know your students as individuals''), not merely aspirational (``be positive'', ``be kind'')? If YES $\rightarrow$ W=4. If NO $\rightarrow$ W=3.

Your reasons must cite specific features from the policy text and explain why the chosen score is not one point higher or lower. CRITICAL: You may ONLY cite text that appears in the policy\_text field of this prompt. Do NOT reproduce or quote text from the rubric anchor descriptions as if it were evidence from the policy being scored. Anchor descriptions are calibration references --- citing them as if they appear in the current policy is fabrication and will produce wrong scores. Return only schema-valid JSON.
\end{quote}

The payload supplies the URN, school name and the title of the selected
document; a set of keyword-matched passages, explicitly labelled as
preliminary signals to be checked against the full text rather than scored on
directly; a flag and term list recording whether specialist-support vocabulary
is present; and the extracted policy text. The required output is a score, a
confidence level and a reason for each of strictness and warmth, plus a
manual-review flag.

\paragraph{Strictness scale anchors}
\begin{description}
\item[Score 1] No consequence language of any kind. The policy is empty, consists entirely of rewards and recognition content, or is so aspirational that no sanctions, detentions, or consequences are described or even implied. Purely aspirational statements about values, culture, or the importance of good behaviour are S=1 if unaccompanied by any consequence content whatsoever.
\item[Score 2] Consequence language is present but no coherent day-to-day system is described. Three common forms: (a) proportionate or contextual sanctions are mentioned without describing when, how, or by whom they are applied; (b) a restorative or relational approach is the primary framework and sanctions appear only as a last resort; (c) the only consequence content is an exclusions/suspensions procedure with no working day-to-day system described alongside it. Key test: could a new member of staff read this policy and know what to do tomorrow when a student disrupts a lesson? If not, it is S=2 (assuming some consequence language is present --- otherwise S=1).
\item[Score 3] A working detentions and isolation system with a clear rules framework. Staff apply consequences by judgment from a menu of options --- no named tracking system, no automatic thresholds. Language such as ``sanctions may include'' and ``at the discretion of staff'' throughout. This is the standard English secondary behaviour policy.
\item[Score 4] A named tracking system where consequence type is determined automatically by position in the system, removing staff discretion over WHICH consequence follows. Two qualifying forms: (a) POINT-ACCUMULATION --- named demerit or BfL points reach a numeric threshold that automatically triggers a named consequence (e.g. ``4 demerits $\rightarrow$ 60-minute detention''); (b) TIERED EVENT-ESCALATION --- named consequence tiers (e.g. C1/C2/C3, Level 1/2/3) where each tier automatically produces a defined sanction (e.g. 'C1 = verbal warning recorded; C2 = same-day detention; C3 = removal to [named facility]'). Form (b) does NOT require a numeric point total --- the tier itself is the decision mechanism. Key test: once a threshold is crossed or a tier is invoked, is the consequence determined by the system position, or by staff judgment at that moment? If the former, S=4. If staff are still choosing from a menu of possible consequences at the point of issue, S=3. Staff may retain discretion over whether to initiate the first action, but not over which consequence the system then generates. Skipping to a higher tier for a more serious incident is severity calibration within S=4, not evidence of discretion over whether to use the system.
\item[Score 5] S=4 baseline, plus two additional features that must BOTH be present: (1) MANDATORY issuance --- the staff member has no discretion at the trigger point (``a demerit is given if a pupil does X'', not ``staff may give a demerit''); AND (2) triggers include continuous IN-LESSON DEPORTMENT NORMS --- always-on conduct standards actively monitored throughout lessons, such as posture/SLANT, eye contact/tracking the teacher, or corridor deportment. WHAT DOES NOT QUALIFY as a deportment-norm trigger for criterion (2): --- Demerit CATEGORY LABELS such as ``Line up conduct'', ``Uniform'', ``Equipment'', ``Transition conduct'' name violation types, not conduct surveillance standards. A list of category labels without accompanying mandatory-issuance language is a category taxonomy --- S=4 maximum. --- Classroom procedural transitions (standing behind chair, lining up) are classroom procedures, not ambient conduct norms --- they do not qualify. --- A stepped-warning process (e.g. reminder $\rightarrow$ verbal warning $\rightarrow$ demerit at step 3) makes issuance discretionary, capping at S=4 regardless of what the category list contains. Score S=5 ONLY when both mandatory issuance AND genuine deportment-norm triggers (posture, SLANT, eye contact) are unambiguously present in the text.
\end{description}

\paragraph{Warmth scale anchors}
\begin{description}
\item[Score 1] No described positive culture infrastructure. The policy may contain an aspirational sentence about values or mention that good behaviour earns a reward, but no reward system is named, no pastoral structures are described, and no relational language about how staff and students interact is present.
\item[Score 2] Aspirational values or ethos language is present (dignity, respect, kindness) but the policy is operationally consequence-dominated. Positive elements exist in name only: rewards mentioned but not described, pastoral support referenced but not structured, relational language generic and aspirational rather than operational.
\item[Score 3] A described rewards system (house points, certificates, commendations, positive contact home) and reference to pastoral support structures. The positive culture section is functional --- processes and systems are named --- but staff-student relational quality is not characterised. This is the standard English secondary policy warmth level.
\item[Score 4] Both of the following clearly present --- these two criteria are necessary AND sufficient. Nothing else is required: (1) a named, described reward system with occasions or criteria; (2) explicit operational guidance on HOW staff should relate to individual students --- instructional rather than aspirational (e.g. ``learn their names'', ``use humour to build bridges'', ``know your students as individuals'', de-escalation scripts, greeting strategies). SCORING RULE: If BOTH criteria are clearly present, assign W=4. Do not require any third criterion. IMPORTANT --- WHAT IS IRRELEVANT TO THIS DECISION: Pastoral support structures (Heads of Year, form tutors, wellbeing teams, safeguarding leads) appear in the W=3 descriptor but play NO role in the W=4 decision. Do not use their presence or absence to determine this score. A school with criterion (1) AND criterion (2) scores W=4 even if no pastoral structures are described.
\item[Score 5] W=4 criteria met, plus the policy explicitly characterises the relational culture --- naming warmth, care, or belonging as values and describing the quality of the staff-student relationship itself, not just what staff should do operationally. This goes beyond tips or instructions into describing how staff and students actually relate (e.g. ``staff know and invest in pupils as individuals'', ``relationships of care and trust underpin how we respond to behaviour''). The positive culture content is substantial and prominently placed in the document, not confined to a brief introduction or appendix.
\end{description}

\paragraph{Scoring rules}
\begin{enumerate}
\item SCORE INDEPENDENTLY: Strictness and warmth are not opposites --- a policy can be high on both or low on both. Score each dimension on its own evidence.
\item MENTION $\neq$ DESCRIBED SYSTEM: The words ``detention'', ``sanction'', or ``reward'' in the text do not mean a described system exists. A described sanctions system explains when consequences are applied, by whom, and the procedure. A described reward system names what rewards exist and how they are given. Do not default to S=3 or W=3 simply because the words appear.
\item EXCLUSION PROCEDURES ARE NOT A SANCTIONS SYSTEM: A policy whose main body describes what offences lead to suspension or permanent exclusion is not a day-to-day behaviour management system. Detailed exclusion procedure text --- even spanning multiple pages --- is consistent with S=2 if no working day-to-day consequence system is described alongside it.
\item S=3 IS NOT THE UNCERTAINTY DEFAULT: S=3 must be earned with a described working detention system. If your reasoning concludes no day-to-day system is described, score S=1 or S=2 --- not S=3 because the evidence is thin.
\item S=4/5 BOUNDARY --- DISCRETION AT THE TRIGGER: S=4 removes discretion over WHICH consequence follows a threshold or tier. S=5 additionally requires that initiating a consequence is MANDATORY at the trigger point (not a choice the staff member makes), AND that triggers include continuous in-lesson deportment norms. Standard rule compliance violations listed as demerit categories --- uniform infringement, equipment, lateness --- do NOT qualify as deportment-norm triggers for S=5. They must be part of an always-on conduct surveillance system explicitly enforced alongside posture/tracking norms to qualify. Classroom procedural transitions (standing behind chair, lining up) do not qualify. Language test: ``a demerit is given if a pupil [does X]'' is mandatory (S=5 possible); ``staff may give a demerit if a pupil [does X]'' is discretionary (S=4 maximum). A plain list of demerit category labels (e.g. ``L1 - Demerit - Uniform''; ``L1 - Demerit - Transition conduct'') without accompanying mandatory-issuance language is a category taxonomy, not a mandatory trigger --- treat as discretionary (S=4 maximum). Similarly, a stepped-warning process (e.g. reminder $\rightarrow$ verbal warning $\rightarrow$ demerit at step 3) means initiation is discretionary even if the demerit category list includes conduct-related items. S=4 consequences are for rule violations (lateness, equipment, uniform, disruption); S=5 additionally covers ambient conduct norms monitored continuously in lessons --- posture, SLANT, eye contact, tracking the teacher. Demerit category labels such as ``Line up conduct'', ``Uniform'', or ``Equipment'' name violation types, not conduct surveillance standards.
\item W=4 REQUIRES BOTH AND ONLY BOTH: The two W=4 criteria (described reward system; operational staff guidance on relating to students as individuals) are jointly necessary and sufficient. If both are clearly present, score W=4. If either is clearly absent, the maximum is W=3. Do NOT treat pastoral support structures as a prerequisite for W=4 --- they appear in the W=3 descriptor but are irrelevant to the W=4 decision. A school with criterion (1) and criterion (2) clearly met scores W=4 even if no independent pastoral structures are described.
\item WARMTH EXCLUSIONS: SEND provision, mental health promotion programmes, safeguarding, and therapeutic interventions do not count toward warmth.
\item SCORE PRIMARY EMPHASIS: Do not reward incidental language. Score what the policy devotes its attention to --- restorative language buried in a consequence-heavy policy does not raise the warmth score; a brief positive section in an otherwise sanctions-focused document does not raise the score.
\item DO NOT REWARD LENGTH: A longer policy does not earn a higher score. Score on the intensity, specificity, and centrality of the evidence.
\item Use only the supplied evidence. Do not infer beyond the text.
\item BORDERLINE $\rightarrow$ LOWER SCORE: When evidence could support two adjacent scores, assign the lower. The distinguishing feature of the higher score must be clearly and unambiguously present.
\item EXPLAIN ADJACENT SCORES: For each score, briefly explain why it is not one point higher and not one point lower.
\end{enumerate}
